\documentclass[10pt,a4paper]{amsart}
\usepackage[margin=19mm]{geometry}
\usepackage{amsmath,amssymb,mathtools}
\usepackage[scr=rsfs,cal=euler]{mathalfa}
\usepackage{xfrac}
\usepackage[unicode,hidelinks,bookmarks=false]{hyperref}
\newcommand{\ie}{i.e.\ }
\newcommand{\cf}{cf.\ }
\newcommand{\resp}{respectively\ }
\newcommand{\Subsection}{\subsection}
\providecommand{\nobreakdash}{\nobreak\hbox{-}}
\setlength{\emergencystretch}{2em}
\multlinegap0pt
\usepackage{amssymb}
\usepackage{xcolor}
\usepackage{tikz}
\newtheorem{proposition}{Proposition}
\title{Symmetric Jacobi Polynomials on a Triangle and Their Spectral Algebra}
\author{Misael E. Marriaga, Miguel A. Pi\~nar}
\date{Source: arXiv:2607.22751v1 (CC BY 4.0)}
\begin{document}
\maketitle

\noindent\emph{Source and licence.} Symmetric Jacobi Polynomials on a Triangle and Their Spectral Algebra. Version arXiv:2607.22751v1 (CC BY 4.0), distributed by arXiv under the Creative Commons Attribution 4.0 International licence, http://creativecommons.org/licenses/by/4.0/, as recorded in the arXiv licence metadata for this exact version and shown on the abstract page https://arxiv.org/abs/2607.22751v1. Full licence notice: \texttt{licenses/arxiv-2607.22751-CC-BY-4.0.txt}.\par\smallskip

\noindent\textit{Editorial scope.} Counted unit: the article's proposition lem:highest\_order\_generated (source0.tex lines 1733--1749). The article's complete original proof of it is delivered with it (source0.tex lines 1751--2002, one \texttt{proof} environment of 252 lines). No mathematics is written, repaired or re-derived: the statement and the proof are the article's own lines, in the article's own notation, and no retained line is substituted or re-flowed.  No further source-local context is needed: every cross-reference of every retained line resolves inside the two delivered spans.  Nothing was omitted from any retained span, so the package carries no omission record.  No retained line contains a citation, so the wrapper carries no bibliography and no bibliography entry.  No retained line contains a figure environment, an image-inclusion command or drawing code, so no image is delivered and the package declares no asset dependency.  Editorial formatting only, in addition: an AMS \texttt{amsart} preamble that restates the article's own declarations and macros used by the retained lines, namely the environments \texttt{proposition} on separate counters, together with the standard packages those lines need.  The retained environments print in this document as Proposition 1 in order of appearance, whereas the article prints the same items with the numbering of its own full document.  The complete native source of the article as distributed in the arXiv e-print for this exact version is bundled unchanged as \texttt{sources/205969/source0.tex}.
\begin{proposition}\label{lem:highest_order_generated}
Let \(\mathcal D\in\mathfrak D_{\alpha,\gamma,\kappa}\) be a nonzero
differential operator. Then \(\mathcal D\) has even order \(2m\).
Moreover, in the variables \((u,t)\), the operator has an expansion of the
form
\begin{equation}\label{eq:highest-order-expansion-generated}
\mathcal D
=
\sum_{\ell=0}^{m}
c_\ell\,
\left(\mathcal D_{1}^{\alpha,\gamma,\kappa}\right)^{m-\ell}
\left(\mathcal D_{2}^{\alpha,\kappa}\right)^{\ell}
+
\textup{terms of order lower than }2m
\end{equation}
with constant coefficients \(c_\ell\).
\end{proposition}

\begin{proof}
Let \(N\) be the order of \(\mathcal D\), and write
\[
\mathcal D
=
\sum_{r+s=N}
b_{r,s}(u,t)\,
\partial_u^{\,r}\partial_t^{\,s}
+
\textup{terms of lower order}.
\]
Since the operators in \(\mathfrak D_{\alpha,\gamma,\kappa}\) have the
same polynomial eigenbasis, \(\mathcal D\) commutes with
\(\mathcal D_1^{\alpha,\gamma,\kappa}\) and
\(\mathcal D_2^{\alpha,\kappa}\). We first determine the terms of highest
order in \(\mathcal D\) from these two commutation relations.

In the variables \((u,t)\), the order-two parts of the two basic
operators are
\begin{align*}
\mathcal D_2^{\alpha,\kappa}
&=
-4t(1-t)\partial_t^2
+
\textup{terms of lower order},\\
\mathcal D_1^{\alpha,\gamma,\kappa}
&=
u(1-u)\partial_u^2
+
\frac{4t(1-t)}{u}\partial_t^2
+
\textup{terms of lower order}.
\end{align*}
Comparing the terms of order \(N+1\) in
\[
[\mathcal D,\mathcal D_2^{\alpha,\kappa}]=0
\]
gives, for \(r+s=N\),
\[
s(1-2t)b_{r,s}(u,t)
-
2t(1-t)\partial_t b_{r,s}(u,t)
=
0.
\]
Hence
\[
b_{r,s}(u,t)
=
d_{r,s}(u)[t(1-t)]^{s/2}.
\]
Since the coefficients of \(\mathcal D\), written in the variables
\((u,t)\), are rational functions of \(u\) and \(t\), the terms with
\(s\) odd must vanish. Thus the terms of order \(N\) in \(\mathcal D\)
can be written as
\begin{equation}\label{eq:D-highest-order-N}
\sum_{j=0}^{\lfloor N/2\rfloor}
d_j(u)
\left(-4t(1-t)\partial_t^2\right)^j
\partial_u^{\,N-2j}.
\end{equation}

The second commutation relation determines the dependence of the
coefficients \(d_j\) on \(u\). Comparing the terms of order \(N+1\) in
\[
[\mathcal D,\mathcal D_1^{\alpha,\gamma,\kappa}]=0
\]
gives
\begin{equation}\label{eq:highest-order-recurrence-compact-N}
\left(
\frac{d_j(u)}
{[u(1-u)]^{\frac{N}{2}-j}}
\right)'
=
\left(\frac{N}{2}-j+1\right)u^{-2}
\frac{d_{j-1}(u)}
{[u(1-u)]^{\frac{N}{2}-j+1}},
\quad
0\leq j\leq \left\lfloor N/2\right\rfloor,
\end{equation}
with the convention \(d_{-1}=0\).

The recurrence also forces \(N\) to be even. For \(j=0\),
\eqref{eq:highest-order-recurrence-compact-N} gives
\[
d_0(u)=C_0[u(1-u)]^{N/2}.
\]
Suppose, to the contrary, that \(N\) is odd. Since \(d_0(u)\) is
rational, the preceding identity forces \(C_0=0\), and hence
\(d_0(u)=0\). Assume inductively that
\[
d_0(u)=\cdots=d_{k-1}(u)=0.
\]
Then the right-hand side of
\eqref{eq:highest-order-recurrence-compact-N} vanishes for \(j=k\), and
therefore
\[
d_k(u)=C_k[u(1-u)]^{N/2-k}.
\]
Since \(N/2-k\) is not an integer, the rationality of \(d_k(u)\) forces
\(C_k=0\). Hence all terms of order \(N\) vanish, contradicting the
definition of \(N\). Thus \(N=2m\).

Set
\[
F_j(u)
=
\frac{d_j(u)}{[u(1-u)]^{m-j}},
\qquad 0\leq j\leq m.
\]
Then \eqref{eq:highest-order-recurrence-compact-N} becomes
\begin{equation}\label{eq:Fj_recurrence}
F_j'(u)
=
(m-j+1)u^{-2}F_{j-1}(u),
\qquad F_{-1}=0.
\end{equation}

It remains to express the functions \(F_j\) as a constant linear
combination of the coefficient families arising from the terms of order
\(2m\) in
\[
\left(\mathcal D_1^{\alpha,\gamma,\kappa}\right)^{m-q}
\left(\mathcal D_2^{\alpha,\kappa}\right)^q,
\qquad 0\leq q\leq m.
\]
For fixed \(q\), these terms can be written as
\begin{equation}\label{eq:Dq-highest-order-2m}
\sum_{j=0}^{m}
p_{j,q}(u)
[u(1-u)]^{m-j}
\left(-4t(1-t)\partial_t^2\right)^j
\partial_u^{\,2m-2j},
\end{equation}
where
\[
p_{j,q}(u)
=
\begin{cases}
0, & j<q,\\[4pt]
(-1)^{j-q}
\binom{m-q}{j-q}
u^{-(j-q)}, & j\ge q.
\end{cases}
\]
The functions \(p_{j,q}\) satisfy
\begin{equation}\label{eq:pjq-recurrence}
p_{j,q}'(u)
=
(m-j+1)u^{-2}p_{j-1,q}(u),
\end{equation}
with the convention \(p_{-1,q}=0\). Indeed, if \(j<q\), then both sides
of \eqref{eq:pjq-recurrence} vanish. If \(j=q\), then \(p_{q,q}=1\) and
\(p_{q-1,q}=0\), so the identity also holds. Finally, if \(j>q\), then
\[
\begin{aligned}
p_{j,q}'(u)
&=
(-1)^{j-q+1}(j-q)
\binom{m-q}{j-q}
u^{-(j-q+1)}
\\
&=
(-1)^{j-q-1}(m-j+1)
\binom{m-q}{j-q-1}
u^{-(j-q+1)}
\\
&=
(m-j+1)u^{-2}p_{j-1,q}(u),
\end{aligned}
\]
where the middle equality uses
\[
(j-q)\binom{m-q}{j-q}
=
(m-j+1)\binom{m-q}{j-q-1}.
\]
Thus, for each \(q\), the functions \(p_{j,q}\) satisfy the same
recurrence as the functions \(F_j\), with \(p_{q,q}=1\) and
\(p_{j,q}=0\) for \(j<q\).

By \eqref{eq:D-highest-order-N}, with \(N=2m\), and
\eqref{eq:Dq-highest-order-2m}, equality of the terms of order \(2m\) is
equivalent to
\[
d_j(u)
=
[u(1-u)]^{m-j}
\sum_{q=0}^{j}c_qp_{j,q}(u),
\qquad 0\le j\le m,
\]
where the sum terminates at \(q=j\) because \(p_{j,q}=0\) for \(q>j\).
In terms of \(F_j\), this becomes the triangular system
\[
F_j(u)
=
\sum_{q=0}^{j}c_qp_{j,q}(u),
\qquad 0\le j\le m.
\]

We solve this system by induction on \(j\). Since \(F_0'=0\) and
\(p_{0,0}=1\), set
\[
c_0=F_0.
\]
Assume that \(c_0,\ldots,c_{j-1}\) have been chosen so that
\[
F_k(u)
=
\sum_{q=0}^{k}c_qp_{k,q}(u),
\qquad 0\le k<j.
\]
Define
\[
H_j(u)
=
F_j(u)
-
\sum_{q=0}^{j-1}c_qp_{j,q}(u).
\]
Applying \eqref{eq:Fj_recurrence} and \eqref{eq:pjq-recurrence} gives
\[
H_j'(u)
=
(m-j+1)u^{-2}
\left(
F_{j-1}(u)
-
\sum_{q=0}^{j-1}c_qp_{j-1,q}(u)
\right).
\]
The expression in parentheses vanishes by the induction hypothesis.
Hence \(H_j\) is constant, and we set
\[
c_j=H_j.
\]
Since \(p_{j,j}=1\), this gives the required identity at index \(j\) and
completes the induction.

Multiplying the resulting identities by \([u(1-u)]^{m-j}\), the
coefficients of the terms of order \(2m\) agree for every
\(0\leq j\leq m\). Therefore
\[
\mathcal D
-
\sum_{q=0}^{m}
c_q
\left(\mathcal D_1^{\alpha,\gamma,\kappa}\right)^{m-q}
\left(\mathcal D_2^{\alpha,\kappa}\right)^q
\]
has order lower than \(2m\).
\end{proof}

\end{document}
